\section{Trabajos Relacionados}
La inspección automatizada de defectos superficiales y clasificación de texturas en materiales de caucho cuenta con una trayectoria significativa en la literatura científica. Los enfoques han evolucionado desde métodos clásicos de procesamiento digital de imágenes hasta el modelado profundo con arquitecturas convolucionales avanzadas.

Históricamente, la detección de anomalías en texturas de llantas se realizaba mediante algoritmos basados en extracción manual de características de baja frecuencia. Trabajos iniciales exploraron el uso de Filtros de Gabor y Patrones Binarios Locales (LBP) para representar la rugosidad y orientación de las estrías. Aunque estos descriptores clásicos eran rápidos de computar, carecían de robustez ante variaciones bruscas de iluminación, sombras complejas causadas por la profundidad de la huella y la presencia de suciedad o agua sobre el caucho.

La revolución de las Redes Neuronales Convolucionales (CNN) profundas, consolidada por \citet{krizhevsky2012imagenet} con el desarrollo de AlexNet, demostró la capacidad de los modelos neuronales para aprender representaciones espaciales jerárquicas directamente a partir de pixeles crudos, superando las limitaciones de los descriptores manuales tradicionales. Posteriormente, \citet{simonyan2014very} exploraron el impacto de la profundidad de la red con la arquitectura VGG, evidenciando que filtros pequeños apilados ($3 \times 3$) capturan mejor las texturas no lineales complejas, lo cual es de gran relevancia para patrones de desgaste en cauchos.

Sin embargo, las redes extremadamente profundas sufrían de degradación del gradiente durante la retropropagación, problema resuelto por \citet{he2016deep} mediante la introducción de conexiones residuales en la arquitectura \texttt{ResNet-50}. Esta red se ha consolidado como un estándar de la industria para tareas de transferencia de aprendizaje (\textit{Transfer Learning}), permitiendo transferir filtros genéricos entrenados en millones de imágenes a conjuntos de datos industriales pequeños con alta estabilidad y convergencia acelerada.

Por otra parte, el problema de la detección de objetos y clasificación ante la presencia de un desbalance de clases severo (frecuente en fallos industriales) fue abordado por \citet{lin2017focal} mediante la formulación de \textit{Focal Loss}. Este avance redefinió la optimización matemática al atenuar dinámicamente la influencia de ejemplos fáciles de clasificar sobre el gradiente, concentrando los pesos en ejemplos difíciles. 

Finalmente, para dotar de interpretabilidad y confianza a los modelos profundo en aplicaciones críticas de inspección visual, \citet{selvaraju2017grad} propusieron Grad-CAM. Esta técnica genera mapas de calor de activación basados en los gradientes de la última capa convolucional, permitiendo verificar visualmente si las decisiones de clasificación se fundamentan en anomalías reales del objeto o en sesgos espaciales indeseados del fondo.

En la actualidad, diversos trabajos aplican estas herramientas de manera conjunta para el análisis industrial de neumáticos. Por ejemplo, sistemas automatizados combinan redes tipo ResNet para la clasificación de defectos en imágenes de rayos X y luz estructurada. Este trabajo se inspira en esta corriente de investigación aplicada, consolidando un estudio de ablación detallado sobre imágenes visuales de textura con un enfoque centrado en mitigar el desbalance textural y proveer interpretabilidad.
